% Ad Atticum 14.15 (ad-atticum-14-15)
% Source: https://raw.githubusercontent.com/PerseusDL/canonical-latinLit/master/data/phi0474/phi057/phi0474.phi057.perseus-lat1.xml
% Fetched by scripts/fetch_latin.py

% Dateline (Perseus): Scr. in Cumano K. Mart. a. 710 (44) .

\ciceroLetterOpener{CICERO ATTICO salutem}

\ciceroSection{1} o mirificum Dolabellam meum! iam enim dico meum; antea, crede mihi, subdubitabam. magnam ἀναθεώρησιν res habet, de saxo, in crucem, columnam tollere, locum illum sternendum locare! quid quaeris? heroica. sustulisse mihi videtur simulationem desideri adhuc quae serpebat in dies, et inveterata verebar ne periculosa nostris tyrannoctonis esset.

\ciceroSection{2} nunc prorsus adsentior tuis litteris speroque meliora. quamquam istos ferre non possum qui, dum se pacem velle simulant, acta nefaria defendunt. sed non possunt omnia simul. incipit res melius ire quam putaram. nec vero discedam nisi cum tu me id honeste putabis facere posse. Bruto certe meo nullo loco deero idque, etiam si mihi cum illo nihil fuisset, facerem propter eius singularem incredibilemque virtutem.

\ciceroSection{3} Piliae nostrae villam totam quaeque in villa sunt trado in Pompeianum ipse proficiscens K. Mai. quam velim Bruto persuadeas ut Asturae sit!
